% PENDING integration/root audit. Written during checkpoint08 freeze.
\begin{saresult}{A.124}{Proved under stated hypotheses}{The component presentation in dimension five when the second homotopy group vanishes}
Let $M$ be a compact connected PL manifold of dimension five,
possibly with boundary, and assume $\pi_2(M)=0$.
For every finite connected pointed homotopy model $K$ of $M$,
\[
 \pi_0P_{\rm PL}(M)\cong E(K),
\]
where $E(K)$ is the complete countable group presentation in
A.119. The same calculation holds for the stipulated definable
concordance components of any definable realization over a real
closed field. It determines the full group, including its
extension law; it is not a finite invariant-factor algorithm.
There is no assertion here for five-manifolds with nonzero
$\pi_2$, or for higher homotopy groups of the comparison.
\end{saresult}
\begin{proof}
First give $M$ a compatible smoothing; its existence will be
verified below. Put $\pi=\pi_1M$ and
\[
 A=\operatorname{Wh}_1^+(\pi;\mathbb Z/2)
   =H_0(\pi;(\mathbb Z/2)[\pi])/(\mathbb Z/2)[1],
 \qquad
 I=\operatorname{im}\bigl(
 \chi_{\rm Wh}:\operatorname{Wh}_3(\pi)\longrightarrow A\bigr).
\]
We invoke the full geometric exact-sequence theorem
\cite[Theorem~7.1, p.~23]{Igusa2025}, which explicitly assumes
$\dim M\geq5$ and $\pi_2M=0$ and includes compact manifolds
with boundary, with the bottom and side fixed. It gives
\[
 0\longrightarrow A/I\xrightarrow{\theta_M}
 \pi_0P_{\rm Diff}(M)\xrightarrow{\Sigma_M}
 \operatorname{Wh}_2(\pi)\longrightarrow0.              \tag{A.124.1}
\]
Exactly the same theorem applies to $M\times I$, after rounding
corners. Projection identifies its fundamental group with
$\pi$, and its second homotopy group is still zero. The map
$\chi_{\rm Wh}$ and hence $I$ are the identical group-theoretic
map and subgroup in these two sequences; naturality is proved
in \cite[Section~4, pp.~14--16]{Igusa2025}.
We use Theorem~7.1 as a cited full theorem. The paper explicitly
describes its geometric proof as partial and refers to the
corrected obstruction theory for the remaining details; we do
not infer this full-kernel assertion from a partial coefficient
calculation.
Two typographical errors in that source are kept separate from
its theorem: in Theorem~7.1(a), p.~23, the printed
$H_2\operatorname{GL}(\mathbb Z[\pi])'\cong K_3(\mathbb Z[\pi])$
has $K_3$ in place of $K_2$; the introductory paragraph on
p.~2 says $\pi_1P$ where its displayed sequence and Theorem~7.1
correctly use $\pi_0P$. We use the standard
$\operatorname{Wh}_2$ and the explicitly displayed component
sequence, consistently with A.119 and A.122.

Consider the actual positive suspension
\[
 \sigma^+:P_{\rm Diff}(M)\longrightarrow
                   P_{\rm Diff}(M\times I).
\]
As in A.122, first normalize concordances on fixed boundary
collars. On the bent product, positive suspension applies the
original concordance on each fixed polar ray. It therefore
preserves composition; the collar normalization is an
equivalence on components. This gives a group homomorphism
$\sigma^+_*$, with the bottom and the entire new side boundary
fixed. The construction is described in
\cite[Section~4.1, p.~8]{Igusa2021Dimension3}.

We check both obstruction maps, rather than identifying only
the abstract end groups in (A.124.1).
The first obstruction is unchanged by suspension:
\[
 \Sigma_{M\times I}\sigma^+_* =\Sigma_M.               \tag{A.124.2}
\]
This is stated in \cite[Part~I, p.~130, following
Section~I.5(5.2), p.~60]{HatcherWagoner1973}; adding a positive
quadratic coordinate leaves the ordered handle-slide matrices
and their labels unchanged.

For the second obstruction it suffices to consider
$\ker\Sigma_M$. Use a one-parameter Morse representative in
the fixed adjacent indices $2,3$. The opened
\cite[Part~I, Chapter~V, Theorem~3.1(a), p.~184;
Part~II, Section~2, p.~244]{HatcherWagoner1973}
permits $2\leq i\leq n-2$, so $i=2$ is permitted when $n=5$.
The corrected geometric invariant in
\cite[Section~7, items~(1)--(12), p.~24]{Igusa2025}
is obtained as follows. The labelled intersection graph $J$
is cut and reassembled using the inverse handle-slide matrix;
for each nonidentity label $u\in\pi$ this gives a closed
framed one-manifold $J_u$. Its framed bordism class is an
element of $\Omega_1^{\rm fr}=\mathbb Z/2$. Summing these
classes with their labels and then passing to conjugacy
coinvariants and to $A/I$ is the inverse of $\theta_M$ on
$\ker\Sigma_M$ in (A.124.1).

In the positive quadratic suspension, the intersections,
their paths to the basepoint, their labels, and the
handle-slide matrices are unchanged. Each $J_u$ is the same
one-manifold in the equatorial copy of the old intersection
space; its normal framing has one additional trivial normal
coordinate. The framed bordism class is unchanged by this
stabilization. Any sign caused by a convention on ordering
normal vectors is also invisible in $\mathbb Z/2$.
The prescribed agreement of framings at birth-death and
handle-addition points is preserved by appending the same
constant normal vector. Thus the entire construction, including
the corner smoothing of $J_u$, commutes with suspension.
Consequently
\[
 \sigma^+_*\theta_M=\theta_{M\times I}.                \tag{A.124.3}
\]
This step uses only framed one-manifolds and their stable
framing classes. It does not invoke the $i,n-i\geq3$
realization lemma for an arbitrary $\pi_2M$ coefficient;
that coefficient is zero under the present hypothesis.
The realization and full kernel needed here are exactly those
already supplied by the cited dimension-five Theorem~7.1.

Equations (A.124.2)--(A.124.3) give a commutative diagram of
the two short exact sequences (A.124.1), with identity maps
on $A/I$ and $\operatorname{Wh}_2(\pi)$. The middle map is
therefore an isomorphism of groups. Explicitly, a suspended
zero first has zero first obstruction, lies in the image of
$\theta_M$, and has zero coefficient there, proving
injectivity. For surjectivity, lift the first obstruction of
a target element to a source element; their difference lies
in the target $\theta$ image, which is also a suspended image.
No splitting of (A.124.1) has been used.

The target $M\times I$ is six-dimensional, so A.122 identifies
its component group with the stable group $E(K)$, since it
has the same homotopy model as $M$. The preceding group
isomorphism proves the smooth dimension-five formula and
preserves the extension law in the countable presentation.

For a PL five-manifold choose a finite PL handle decomposition,
including collars when its boundary is nonempty. Extend a
compatible smoothing over handles in increasing index.
The classical PL--DIFF handle argument in
\cite[Essay~IV, Appendix~B, p.~211]{KirbySiebenmann1977}
places the index-$k$ obstruction in $\Gamma_{k-1}$.
All indices are at most five. The low-dimensional compatible
smoothings in \cite[Essay~I, p.~39]{KirbySiebenmann1977}
and the explicit equality $\Gamma_4=0$ in
\cite[Essay~IV, p.~200]{KirbySiebenmann1977}
therefore remove every obstruction. This constructs a
compatible smoothing of the given PL manifold; no smoothing
has been prescribed independently on its boundary.
The smooth-to-PL component comparison in A.103, valid in
dimension at least five and relative to the side boundary,
transfers the smooth calculation. Finally A.113 and A.115,
with the same relative boundary convention, transfer the
component group to every definable realization and every
real closed field. The finite model of $M$ in A.119 may be
replaced by any finite pointed homotopy model by the
functorial group identification proved there.
\end{proof}
